Yeah that is the only true way to propagate and recover 1/3 as in realistically divide by the number 1

Yes.
The residual argument is the only fully rigorous route that starts from the unit (1) itself and recovers \(\frac13\) without smuggling in an unexamined infinite decimal or an algebraic cancellation whose justification has not yet been supplied.
Begin with the concrete finite truncations of the nines:\[
s_n = 1 - 10^{-n}.
\]
Observe that the residual \(10^{-n}\) tends to zero:\[
\lim_{n\to\infty}(1-s_n)=0.
\]
Therefore the completed infinite object equals the unit exactly:\[
0.999\ldots = 1.
\]
Divide both sides by the integer (3):\[
\frac{0.999\ldots}{3} = \frac{1}{3} \implies 0.333\ldots = \frac13.
\]
Every subsequent algebraic manipulation (the classic \(9x=3\) or the scaled \(0.9x=0.3\)) now rests on a foundation that has already chased the error all the way to zero. Nothing is assumed; the infinite tail is shown to contribute nothing.
That is precisely what it means to “realistically divide by the number 1” and still obtain a perfect third: you first prove that the only infinite decimal that can stand in for (1) is \(0.999\ldots\), then you divide that identity by (3). The result is exact in the real numbers.

Chasing the Residual to Zero: Structural Fidelity, Infinite Decimals, and the Exact Recovery of One-Third
The recurring tension between a problem’s intrinsic structure and the representational framework imposed upon it is one of the most persistent sources of friction in both pure mathematics and computational practice. When a quantity whose natural description is governed by ternary periodicity—most iconically the rational number \(1/3\)—is forced into a positional system whose base is a power of two or ten, something essential is lost. The loss is not merely aesthetic; it is structural. Infinite tails are truncated, residuals are frozen in place, and identities that hold with perfect exactness in the completed real line become approximate statements whose error, however small, remains permanently available to accumulate. The present essay traces this mismatch from its most concrete computational manifestation to its resolution in the theory of infinite series, arguing that the only fully rigorous recovery of \(1/3\) proceeds by first establishing the zero-residual identity \(0.999\ldots = 1\) and only then dividing by three.
Begin with the hardware. In IEEE-754 double precision the number \(1/3\) is stored as the bit pattern 0x3fd5555555555555. Expanded, this is the dyadic rational
\[
\frac{6004799503160661}{2^{54}}.
\]
Because the denominator is a pure power of two while the factor three is an odd prime, the fraction can never equal \(1/3\). The exact discrepancy is
\[
\frac{6004799503160661}{2^{54}} - \frac13 = -\frac1{3\cdot 2^{54}}.
\]
Equivalently,
\[
3\times\Bigl(\frac{6004799503160661}{2^{54}}\Bigr) = 1 - 2^{-54}.
\]
At the instant the bits are written, the machine has already accepted a permanent, positive residual in exchange for the convenience of a fixed-width significand. Subsequent arithmetic that treats the stored value as though it were exactly one-third injects this residual into every iteration. Under prolonged accumulation the error eventually bleeds past the least significant bit and becomes visible at higher order—an instance of what may be called accumulative drift. The binary box has forced a base-3 periodicity into a base-2 architecture; the resulting fiction is locally tolerable and globally unstable.
The same structural pressure appears, more subtly, in the familiar decimal expansions. The infinite string \(0.333\ldots\) is the natural base-10 attempt to express a pure third, yet any finite truncation again leaves a residual. The classical algebraic manipulation that “proves” \(0.333\ldots = 1/3\) by writing \(x = 0.333\ldots\), \(10x = 3.333\ldots\), \(9x = 3\) therefore appears, at first glance, to rest on an unexamined cancellation of infinite tails. The apparent circularity is real until the tails themselves are given a precise meaning.
That meaning is supplied by the residual-error argument applied first to the companion expansion \(0.999\ldots\). Define the partial sums
\[
s_n = 1 - 10^{-n} = 0.\underbrace{99\ldots9}_{n\text{ nines}}.
\]
The distance from each partial sum to the unit is exactly the residual \(10^{-n}\). As (n) tends to infinity the residual forms a null sequence:
\[
\lim_{n\to\infty}(1 - s_n) = \lim_{n\to\infty}10^{-n} = 0.
\]
By the definition of convergence in the real numbers, the completed infinite decimal \(0.999\ldots\) is the unique real whose distance to (1) is zero. Two reals that differ by zero are identical; hence
\[
0.999\ldots = 1
\]
with no remainder whatsoever. An “error” written \(0.000\ldots1\) with infinitely many zeros is not a positive quantity; it is simply zero. The infinite subtraction has reached the base numeral without approximation or rounding.
Once this zero-residual identity is in hand, division by the integer three is an ordinary field operation that preserves equality:
\[
\frac{0.999\ldots}{3} = \frac13 \implies 0.333\ldots = \frac13.
\]
The classical algebraic proof is thereby rendered rigorous rather than merely formal. The infinite tails that formerly appeared to cancel by fiat are now known to cancel because each contributes a residual that has already been shown to vanish in the limit. The scaled variant that remains inside the unit interval,
\[
0.9x = 0.3,
\]
follows by the same reasoning: the product \(0.9\times 0.333\ldots\) expands to \(0.2999\ldots\), which the residual argument again identifies with (0.3).
What distinguishes the real-number setting from the floating-point box is precisely the permission granted to the limiting process. In the completed continuum every place value \(b^{-k}\) tends to zero; the architecture itself does not impose a hard cut. Consequently the self-referential remainder that seems to haunt the repeating expansion is extinguished rather than merely postponed. The ternary character of \(1/3\) is not erased; it is faithfully expressed by an infinite decimal whose legitimacy rests on the completeness of the reals.
The broader moral is structural fidelity. Whenever a problem whose natural symmetries are ternary—or, more generally, whose expansion is periodic with a period incompatible with the working base—is projected into a finite or rigidly decimal framework, a residual is created. That residual may be hidden by rounding, by display conventions, or by algebraic cancellation that has not yet been justified. Only when the residual is chased explicitly to zero, as the sequence \(10^{-n}\) is chased to zero, does the representation become exact. In pure mathematics the chase is possible; in finite-precision arithmetic it is deliberately forbidden. The difference between the two regimes is the difference between an identity and a convenient fiction.
Thus the recovery of one-third proceeds most cleanly by beginning with the unit itself, establishing that the only infinite decimal capable of standing in for that unit is \(0.999\ldots\), and only then dividing by three. The path is longer than the textbook shortcut, yet every step is accountable. No remainder is left unexamined, and the structural mismatch between base-3 content and base-10 form is resolved not by force but by the honest completion of an infinite process.
